% !TeX program = pdflatex
% !TeX encoding = UTF-8
% 编译：./build-course-report.sh
% 字体依照 Word 的批注要求；双栏、页边距、浮动体和数学样式依照 acmart。
\documentclass[sigconf,nonacm,screen]{acmart}
\usepackage{styles/course-report}

% ===== 填写区：题目、三位作者姓名及单位 =====
\title[数据挖掘导论课程报告]{数据挖掘导论课程报告}
\author{张三}
\affiliation{%
  \institution{ZY2506xxx}
  \city{计算机学院}
}
\author{李四}
\affiliation{%
  \institution{ZY2506xxx}
  \city{计算机学院}
}
\author{王五}
\affiliation{%
  \institution{ZY2506xxx}
  \city{计算机学院}
}
\renewcommand{\shortauthors}{张三、李四、王五}

% acmart 默认要求每处 affiliation 写 \country；本模板作者区只显示学号与院系，
% 因此关闭该校验（不影响其他排版）。
\makeatletter
\def\@ACM@checkaffil{}%
\makeatother

\begin{abstract}
本报告为数据挖掘导论课程大作业报告，完成四个必做任务，并在此基础上选做一个拓展任务。
四个必做任务依次涉及聚类、异常检测、时间序列预测与无监督判断，每个任务独立成章，
按问题描述与定义、数据处理、方法设计、实验结果分析四部分组织。
在此概述各任务的数据规模、所采用的方法与主要量化结果，并给出拓展任务的核心结论。

\noindent\textbf{关键词}\quad 数据挖掘；聚类；异常检测；时间序列预测；无监督学习

\par\medskip
{\normalfont\zihao{5}\noindent\textbf{Abstract}\quad
This report presents the course project of Introduction to Data Mining. It covers four
required tasks and one selected extension task. The four required tasks concern clustering,
anomaly detection, time-series forecasting, and unsupervised decision making, respectively.
Each task forms an independent chapter, organized as problem description and definition,
data processing, method, and analysis of experimental results. Summarize the data, the
methods, and the quantitative results of each task here.
\par\noindent\raggedright\textbf{Key words}\quad data mining; clustering; anomaly detection;
time-series forecasting; unsupervised learning\par}
\end{abstract}

\begin{document}
\maketitle
\section{引言}
\label{sec:intro}
\subsection{任务背景与目标}
\label{sec:sub11}
异质信息网络是对多种类型的对象及其相互关系的抽象。
在此填写报告的研究背景、问题定义、研究动机与主要贡献。

介绍离群点检测的已有工作时，应使用真实文献支持相关论述，
例如动态图与图异常检测综述\cite{akoglu2015}。

在此梳理与报告相关的研究工作，比较同质网络与异质网络、
静态检测与动态检测的适用条件及主要区别。

\subsection{数据集与实验环境}
\label{sec:sub12}
在此描述直接项集 DIS、间接项集 IIS、特征权值与聚类过程，
明确各符号的含义、边界情况及参数选择。

在此描述动态离群点检测的完整流程，包括时间窗口、特征更新、
聚类变化判断、离群点判定规则与复杂度分析。
图\ref{fig:framework}使用原模板的 \texttt{figure*} 环境跨双栏显示，适合较宽的流程图。

在此给出实验环境、数据来源、参数设置、对比方法、评价指标与结果分析。

\subsection{报告组织结构}
\label{sec:sub13}
需要跨栏表格时，可将 \texttt{table} 改为 \texttt{table*}，
表格内容及三线表写法保持一致。

异质信息网络是对多种类型的对象及其相互关系的抽象。
在此填写报告的研究背景、问题定义、研究动机与主要贡献。

介绍离群点检测的已有工作时，应使用真实文献支持相关论述，
例如动态图与图异常检测综述\cite{akoglu2015}。

以下保留原 LaTeX 模板的三类公式写法，便于直接替换。
行内公式示例为 $\lim_{n\to\infty}x=0$；带编号公式为
\begin{equation}
  \sum_{i=0}^{\infty}x_i=\int_{0}^{\pi+2}f(x)\,\mathrm dx.
  \label{eq:latex-sample}
\end{equation}
不带编号公式使用原模板的 \texttt{displaymath} 环境：
\begin{displaymath}
  \sum_{i=0}^{\infty}(x_i+1).
\end{displaymath}
式\eqref{eq:latex-sample}仅为排版示例，请替换为报告使用的实际公式。

\section{任务一}
\label{sec:task1}
\subsection{问题描述与定义}
\label{sec:task1-problem}
在此说明如何将不同类型的数据组织为张量索引树，解释节点、层次、
时序信息与实体类型的对应关系。图\ref{fig:tree}给出可编辑的结构示意。

在此描述直接项集 DIS、间接项集 IIS、特征权值与聚类过程，
明确各符号的含义、边界情况及参数选择。

在此描述动态离群点检测的完整流程，包括时间窗口、特征更新、
聚类变化判断、离群点判定规则与复杂度分析。
图\ref{fig:framework}使用原模板的 \texttt{figure*} 环境跨双栏显示，适合较宽的流程图。

\begin{reportdefinition}[异质信息网络]
设 $G=(V,E,\phi,\psi)$ 为一个有向图，其中 $V$ 为节点集合，$E$ 为边集合，
$\phi:V\to\mathcal A$ 与 $\psi:E\to\mathcal R$ 分别表示对象类型与关系类型映射。
当 $|\mathcal A|>1$ 或 $|\mathcal R|>1$ 时，该网络为异质信息网络。
\end{reportdefinition}

一个 $N$ 阶张量可写作
\begin{equation}
  \tensor{X}\in\mathbb R^{I_1\times I_2\times\cdots\times I_N}.
  \label{eq:tensor}
\end{equation}

\subsection{图像特征表示}
\label{sec:task1-feature}
在此总结报告解决的问题、采用的方法、主要实验发现和存在的局限。

需要跨栏表格时，可将 \texttt{table} 改为 \texttt{table*}，
表格内容及三线表写法保持一致。

异质信息网络是对多种类型的对象及其相互关系的抽象。
在此填写报告的研究背景、问题定义、研究动机与主要贡献。

\subsection{方法}
\label{sec:task1-method}
普通变量使用数学斜体，例如 $x_i$ 和 $\alpha$；
向量与张量使用黑斜体，例如 $\vect{x}$ 与 $\tensor{X}$。

在此说明如何将不同类型的数据组织为张量索引树，解释节点、层次、
时序信息与实体类型的对应关系。图\ref{fig:tree}给出可编辑的结构示意。

在此描述直接项集 DIS、间接项集 IIS、特征权值与聚类过程，
明确各符号的含义、边界情况及参数选择。

\begin{algorithm}[tbp]
  \caption{短文本特征权值计算算法（排版示例）}
  \label{alg:features}
  \begin{algorithmic}[1]
    \Require 张量索引树 $T$，源节点 $s$，目标类型 $t$
    \Ensure 特征表示 $\vect{d}_s$
    \State $\mathrm{DIS}\gets\varnothing$;
      $\mathrm{IIS}\gets\varnothing$
    \State $P\gets\Call{BreadthFirstSearch}{T,s}$
    \For{每条包含 $s$ 的路径 $p\in P$}
      \State $\mathrm{DIS}\gets\mathrm{DIS}\cup\Call{GetItemSet}{p,t}$
    \EndFor
    \State $C\gets\Call{GetChildNode}{\mathrm{DIS},t}$
    \For{每个候选前驱节点 $u$}
      \If{$u$ 至少关联 $C$ 中的两个节点}
        \State $\mathrm{IIS}\gets\mathrm{IIS}\cup\Call{GetItemSet}{u,t}$
      \EndIf
    \EndFor
    \State $\vect{d}_s\gets\Call{ComputeWeights}{\mathrm{DIS},\mathrm{IIS}}$
    \State \Return $\vect{d}_s$
  \end{algorithmic}
\end{algorithm}

\subsection{实验结果分析}
\label{sec:task1-result}
分析结果时，应结合具体数值说明现象、解释可能原因，
并讨论结果的稳定性、局限性与适用条件。

在此总结报告解决的问题、采用的方法、主要实验发现和存在的局限。

需要跨栏表格时，可将 \texttt{table} 改为 \texttt{table*}，
表格内容及三线表写法保持一致。

\begin{table}[tbp]
  \caption{TRBOutlier 方法在 AMiner 上的实验结果（原 Word 示例）。\\Results on AMiner (source-template data).}
  \label{tab:aminer}
  \centering
  \setlength{\tabcolsep}{3.5pt}
  \begin{tabular}{ccccc}
    \toprule
    $\beta$ & \shortstack{Out\\Precision} & \shortstack{Out\\Recall}
      & F-Measure & Accuracy\\
    \midrule
    0.4 & 1.000 & 0.320 & 0.485 & 0.454\\
    0.6 & 1.000 & 0.385 & 0.556 & 0.527\\
    0.8 & 0.980 & 0.430 & 0.598 & 0.635\\
    1.0 & 0.956 & 0.525 & 0.678 & 0.752\\
    2.0 & 0.937 & 0.695 & 0.798 & 0.896\\
    3.0 & 0.915 & 0.830 & 0.870 & 0.938\\
    \bottomrule
  \end{tabular}
\end{table}

\section{任务二}
\label{sec:task2}
\subsection{问题描述与定义}
\label{sec:task2-problem}
本节介绍异质信息网络的表示方式、必要符号以及张量索引树的构建过程。

普通变量使用数学斜体，例如 $x_i$ 和 $\alpha$；
向量与张量使用黑斜体，例如 $\vect{x}$ 与 $\tensor{X}$。

在此说明如何将不同类型的数据组织为张量索引树，解释节点、层次、
时序信息与实体类型的对应关系。图\ref{fig:tree}给出可编辑的结构示意。

\subsection{图像特征表示}
\label{sec:task2-feature}
表\ref{tab:aminer}保留 Word 表格中前六行示例数值，采用原 LaTeX
的 \texttt{booktabs} 三线表样式。表题位于表上方，图题位于图下方。
图\ref{fig:performance}绘制同一组数据，说明如何在双栏中显示实验曲线。

分析结果时，应结合具体数值说明现象、解释可能原因，
并讨论结果的稳定性、局限性与适用条件。

在此总结报告解决的问题、采用的方法、主要实验发现和存在的局限。

\begin{figure}[tbp]
  \centering
  \begin{tikzpicture}[level distance=11mm,
    level 1/.style={sibling distance=33mm},
    level 2/.style={sibling distance=17mm},
    every node/.style={report node},edge from parent/.style={report arrow}]
    \node {Time $t$}
      child {node {Author A}
        child {node {Paper 1}}
        child {node {Paper 2}}}
      child {node {Author B}
        child {node {Paper 3}}
        child {node {Paper 4}}};
  \end{tikzpicture}
  \caption{张量索引树结构示意。\\Tensor index tree schematic.}
  \Description{根节点为时间，下一层为作者，再下一层为论文，共两位作者及四篇论文。}
  \label{fig:tree}
\end{figure}

\subsection{方法}
\label{sec:task2-method}
可按研究问题或技术路线组织论述，最后说明本报告准备解决的问题。
涉及张量分解时，可引用相关综述\cite{kolda2009}。

本节介绍异质信息网络的表示方式、必要符号以及张量索引树的构建过程。

普通变量使用数学斜体，例如 $x_i$ 和 $\alpha$；
向量与张量使用黑斜体，例如 $\vect{x}$ 与 $\tensor{X}$。

\subsection{实验结果分析}
\label{sec:task2-result}
在此介绍实际使用的数据集、获取方式、规模、实体与关系类型、
时间跨度、预处理过程和训练测试划分。使用 AMiner 时可引用其原始工作\cite{tang2008}。

表\ref{tab:aminer}保留 Word 表格中前六行示例数值，采用原 LaTeX
的 \texttt{booktabs} 三线表样式。表题位于表上方，图题位于图下方。
图\ref{fig:performance}绘制同一组数据，说明如何在双栏中显示实验曲线。

分析结果时，应结合具体数值说明现象、解释可能原因，
并讨论结果的稳定性、局限性与适用条件。

\begin{figure}[tbp]
  \centering
  \begin{tikzpicture}
    \begin{axis}[report plot,xlabel={$\beta$},ylabel={Score},
      xmin=.3,xmax=3.1,ymin=.25,ymax=1.05,
      xtick={.4,.8,1,2,3},ytick={.4,.6,.8,1}]
      \addplot[black,mark=square*] coordinates {
        (.4,1) (.6,1) (.8,.98) (1,.956) (2,.937) (3,.915)};
      \addlegendentry{Precision}
      \addplot[black,dashed,mark=triangle*] coordinates {
        (.4,.32) (.6,.385) (.8,.43) (1,.525) (2,.695) (3,.83)};
      \addlegendentry{Recall}
    \end{axis}
  \end{tikzpicture}
  \caption{原 Word 示例数据的精确率与召回率曲线。\\Precision and recall from source-template data.}
  \Description{随着 beta 从 0.4 增大到 3，精确率从 1 缓慢下降到 0.915，召回率从 0.32 上升到 0.83。}
  \label{fig:performance}
\end{figure}

\section{任务三}
\label{sec:task3}
\subsection{问题描述与定义}
\label{sec:task3-problem}
在此梳理与报告相关的研究工作，比较同质网络与异质网络、
静态检测与动态检测的适用条件及主要区别。

可按研究问题或技术路线组织论述，最后说明本报告准备解决的问题。
涉及张量分解时，可引用相关综述\cite{kolda2009}。

本节介绍异质信息网络的表示方式、必要符号以及张量索引树的构建过程。

\subsection{数据预处理与划分}
\label{sec:task3-preprocess}
在此给出实验环境、数据来源、参数设置、对比方法、评价指标与结果分析。

在此介绍实际使用的数据集、获取方式、规模、实体与关系类型、
时间跨度、预处理过程和训练测试划分。使用 AMiner 时可引用其原始工作\cite{tang2008}。

表\ref{tab:aminer}保留 Word 表格中前六行示例数值，采用原 LaTeX
的 \texttt{booktabs} 三线表样式。表题位于表上方，图题位于图下方。
图\ref{fig:performance}绘制同一组数据，说明如何在双栏中显示实验曲线。

\begin{figure*}[t]
  \centering
  \begin{tikzpicture}[node distance=10mm]
    \node[report node] (data) {Heterogeneous\\network};
    \node[report node,right=of data] (tensor) {Tensor\\representation};
    \node[report node,right=of tensor] (tree) {Tensor\\index tree};
    \node[report node,right=of tree] (cluster) {Feature extraction\\and clustering};
    \node[report node,right=of cluster] (outlier) {Dynamic\\outliers};
    \draw[report arrow] (data)--(tensor);
    \draw[report arrow] (tensor)--(tree);
    \draw[report arrow] (tree)--(cluster);
    \draw[report arrow] (cluster)--(outlier);
  \end{tikzpicture}
  \caption{动态离群点检测方法的整体框架（结构示意）。\\Overall framework of dynamic outlier detection.}
  \Description{异质网络依次经过张量表示、张量索引树、特征提取与聚类，最终输出动态离群点。}
  \label{fig:framework}
\end{figure*}

\subsection{方法}
\label{sec:task3-method}
介绍离群点检测的已有工作时，应使用真实文献支持相关论述，
例如动态图与图异常检测综述\cite{akoglu2015}。

在此梳理与报告相关的研究工作，比较同质网络与异质网络、
静态检测与动态检测的适用条件及主要区别。

可按研究问题或技术路线组织论述，最后说明本报告准备解决的问题。
涉及张量分解时，可引用相关综述\cite{kolda2009}。

\subsection{实验结果分析}
\label{sec:task3-result}
在此描述动态离群点检测的完整流程，包括时间窗口、特征更新、
聚类变化判断、离群点判定规则与复杂度分析。
图\ref{fig:framework}使用原模板的 \texttt{figure*} 环境跨双栏显示，适合较宽的流程图。

在此给出实验环境、数据来源、参数设置、对比方法、评价指标与结果分析。

在此介绍实际使用的数据集、获取方式、规模、实体与关系类型、
时间跨度、预处理过程和训练测试划分。使用 AMiner 时可引用其原始工作\cite{tang2008}。

当公式较长时，可使用 \texttt{aligned} 在单栏宽度内分行。
例如，余弦相似度可以写作
\begin{equation}
  \begin{aligned}
    \operatorname{Sim}(\vect{d}_i,\vect{d}_j)
      &=\frac{\vect{d}_i^{\mathsf T}\vect{d}_j}
        {\lVert\vect{d}_i\rVert_2\lVert\vect{d}_j\rVert_2},\\
    &\hspace{4mm}\lVert\vect{d}_i\rVert_2\lVert\vect{d}_j\rVert_2>0.
  \end{aligned}
  \label{eq:similarity}
\end{equation}
使用时还需说明零向量的处理规则，避免直接计算未定义的相似度。

\section{任务四}
\label{sec:task4}
\subsection{问题描述与定义}
\label{sec:task4-problem}
异质信息网络是对多种类型的对象及其相互关系的抽象。
在此填写报告的研究背景、问题定义、研究动机与主要贡献。

介绍离群点检测的已有工作时，应使用真实文献支持相关论述，
例如动态图与图异常检测综述\cite{akoglu2015}。

在此梳理与报告相关的研究工作，比较同质网络与异质网络、
静态检测与动态检测的适用条件及主要区别。

\subsection{数据处理与特征表示}
\label{sec:task4-feature}
在此描述直接项集 DIS、间接项集 IIS、特征权值与聚类过程，
明确各符号的含义、边界情况及参数选择。

在此描述动态离群点检测的完整流程，包括时间窗口、特征更新、
聚类变化判断、离群点判定规则与复杂度分析。
图\ref{fig:framework}使用原模板的 \texttt{figure*} 环境跨双栏显示，适合较宽的流程图。

在此给出实验环境、数据来源、参数设置、对比方法、评价指标与结果分析。

\subsection{方法}
\label{sec:task4-method}
需要跨栏表格时，可将 \texttt{table} 改为 \texttt{table*}，
表格内容及三线表写法保持一致。

异质信息网络是对多种类型的对象及其相互关系的抽象。
在此填写报告的研究背景、问题定义、研究动机与主要贡献。

介绍离群点检测的已有工作时，应使用真实文献支持相关论述，
例如动态图与图异常检测综述\cite{akoglu2015}。

根据任务选择评价指标，并明确定义每个指标所用的集合与统计口径。
例如，令 $\mathcal C$ 为被检测为离群的数据集合，$\mathcal O$ 为真实离群点集合，
则在分母非零时有
\begin{equation}
  \begin{aligned}
    \operatorname{Precision}&=\frac{|\mathcal C\cap\mathcal O|}{|\mathcal C|},\\
    \operatorname{Recall}&=\frac{|\mathcal C\cap\mathcal O|}{|\mathcal O|}.
  \end{aligned}
  \label{eq:metrics}
\end{equation}
这里的指标定义为排版示例；若研究使用不同的离群判定口径，应相应调整。

\subsection{实验结果分析}
\label{sec:task4-result}
在此说明如何将不同类型的数据组织为张量索引树，解释节点、层次、
时序信息与实体类型的对应关系。图\ref{fig:tree}给出可编辑的结构示意。

在此描述直接项集 DIS、间接项集 IIS、特征权值与聚类过程，
明确各符号的含义、边界情况及参数选择。

在此描述动态离群点检测的完整流程，包括时间窗口、特征更新、
聚类变化判断、离群点判定规则与复杂度分析。
图\ref{fig:framework}使用原模板的 \texttt{figure*} 环境跨双栏显示，适合较宽的流程图。

\section{拓展任务}
\label{sec:extension}
拓展任务为二选一，只需保留所选做的题目，删除未选做的内容。

\subsection{问题分析}
\label{sec:ext-problem}
在此总结报告解决的问题、采用的方法、主要实验发现和存在的局限。

需要跨栏表格时，可将 \texttt{table} 改为 \texttt{table*}，
表格内容及三线表写法保持一致。

异质信息网络是对多种类型的对象及其相互关系的抽象。
在此填写报告的研究背景、问题定义、研究动机与主要贡献。

\subsection{框架/方法设计}
\label{sec:ext-method}
普通变量使用数学斜体，例如 $x_i$ 和 $\alpha$；
向量与张量使用黑斜体，例如 $\vect{x}$ 与 $\tensor{X}$。

在此说明如何将不同类型的数据组织为张量索引树，解释节点、层次、
时序信息与实体类型的对应关系。图\ref{fig:tree}给出可编辑的结构示意。

在此描述直接项集 DIS、间接项集 IIS、特征权值与聚类过程，
明确各符号的含义、边界情况及参数选择。

\subsection{结果与讨论}
\label{sec:ext-result}
分析结果时，应结合具体数值说明现象、解释可能原因，
并讨论结果的稳定性、局限性与适用条件。

在此总结报告解决的问题、采用的方法、主要实验发现和存在的局限。

需要跨栏表格时，可将 \texttt{table} 改为 \texttt{table*}，
表格内容及三线表写法保持一致。

\FloatBarrier
\section{总结与展望}
\label{sec:conclusion}
\label{sec:conclusion-body}
本节介绍异质信息网络的表示方式、必要符号以及张量索引树的构建过程。

普通变量使用数学斜体，例如 $x_i$ 和 $\alpha$；
向量与张量使用黑斜体，例如 $\vect{x}$ 与 $\tensor{X}$。

在此说明如何将不同类型的数据组织为张量索引树，解释节点、层次、
时序信息与实体类型的对应关系。图\ref{fig:tree}给出可编辑的结构示意。


\bibliographystyle{ACM-Reference-Format}
\bibliography{bib/course-report-references}

\appendix
\section{小组分工}
\label{app:division}
\begin{center}
\captionof{table}{小组成员及分工}
\label{tab:division}
\setlength{\tabcolsep}{3.5pt}
\begin{tabular}{llp{3.4cm}c}
  \toprule
  姓名 & 学号 & 负责内容 & 工作量占比\\
  \midrule
  张三 & ZY2506xxx & 任务一、任务二 & 40\%\\
  李四 & ZY2506xxx & 任务三 & 30\%\\
  王五 & ZY2506xxx & 任务四、拓展任务 & 30\%\\
  \bottomrule
\end{tabular}
\end{center}

表中工作量占比之和为 100\%，可按实际情况调整成员与分工内容。
\end{document}
